  \documentclass[twoside=false, %  doppelseitiger Druck
    DIV=15,% DIV Faktor für Satzspiegelberechnung, sie Doku zu KOMA Script
    BCOR=15mm, % Bindekorrektur
    chapterprefix=false,
    headinclude=true,
    footinclude=false,
    pagesize,%         write pagesize to DVI or PDF
    fontsize=11pt,%             use this font size
    paper=a4,%          use ISO A4
    bibliography=totoc,%         write bibliography-chapter to table of contents
    index=totoc,%         write index-chapter to table of contents
    cleardoublepage=plain,% \cleardoublepage generates pages with pagestyle empty
     headings=big,%       A4/B5
    listof=flat,%        improved list of tables
    numbers=noenddot
  ]{scrbook}

\usepackage[utf8]{inputenc}
\usepackage{makeidx}
\usepackage{amsfonts}
\usepackage[slantedGreek,sc]{mathpazo}  % Schriftart Palatino
% \usepackage{lmodern}    % statt mathpazo, falls CM Fonts verwendet werden sollen
%\usepackage{mathptmx}    % statt mathpazo, falls Times  verwendet werden soll
\usepackage[scaled=.95]{helvet}
\usepackage{courier}
\usepackage[T1]{fontenc}
\usepackage{textcomp}
\usepackage{amsmath}            % standard math notation (vectors/sets/...)
\usepackage{bm}        % standard math notation (fonts)
\usepackage{fixmath}        % standard math notation (fonts)
\usepackage{graphicx}
\usepackage[facing=yes]{floatrow}       % mehrere Gleitobjekte nebeneinander/caption neben Bild/Tabelle
\usepackage[labelfont=bf,sf,font=small,labelsep=space,format=plain]{caption}
\usepackage{subcaption}
\usepackage{scrlayer-scrpage}
% \usepackage{pstool}  % einbinden falls psfrag verwendet werden soll
\usepackage{epstopdf}
\usepackage[ngerman]{babel}
\usepackage{ellipsis}  % Korrigiert den Weißraum um Auslassungspunkte
\usepackage{microtype}  % optischer Randausgleich etc.

\usepackage{xcolor}         % z.B. für schattierte Boxen
\usepackage{framed}			% shaded Umgebung
\definecolor{shadecolor}{gray}{.85}%

% Links im PDF
\usepackage[colorlinks=false,
            pdfborder={0 0 0},
            breaklinks=true]
            {hyperref}


%\typearea[current]{calc}


% Einstellungen für Bild-/Tabellenbeschriftung neben dem Bild
\floatsetup[figure]{capbesideposition={inside,top}}
\floatsetup[table]{capbesideposition={inside,top},style=plaintop}
\renewfloatcommand{fcapside}{figure}[\capbeside][\FBwidth]
\newfloatcommand{tcapside}{table}[\capbeside][\FBwidth]


\selectlanguage{ngerman}


\deffootnote{1em}{1em}{%
 \makebox[1em][l]{\thefootnotemark}}

\makeindex

\newcommand{\real}{\mathord{\mathrm{I\!R}}}

\begin{document}
\selectlanguage{ngerman}
\def\figdir{figures}
\def\tabledir{tables}

\frontmatter

\pagestyle{scrplain}
\pagestyle{empty}

\begin{titlepage}

\sffamily

\raggedleft

\vspace*{-2cm}

\includegraphics{\figdir/logo-th-rosenheim-2019_master_quer_2c.eps}

\vfill

\centering
\LARGE
% \vspace*{\fill}
%-----------
Fakultät für Informatik  \vspace{0.5cm}\\
\Large
Studiengang Software- und Systems-Engineering

\vspace{2cm}

\LARGE

Erkennung von Design Patterns in Quellcode durch Machine Learning

\vspace{2cm}

\Large
Master Thesis

\vspace{1.5cm}


\Large
von

\vspace{0.5cm}

%\vspace*{\fill}

\LARGE
Mehmet Aslan \vspace{1cm}

\vspace{1cm}

\flushleft
 \Large
\vspace*{\fill}

%-----------
\begin{tabbing}
Datum der Abgabe: \= 18.02.2024 \kill
Datum der Abgabe: \> 18.02.2024 \\
Erstprüfer: \> Prof.\ Dr.\ Marcel Tilly\\
Zweitprüfer: \> Prof.\ Dr.\ Kai Höfig
\end{tabbing}
%-----------

\end{titlepage}

\cleardoubleemptypage

{
\large
\thispagestyle{empty}
\vspace*{\fill}

\noindent
\textsc{Eigenständigkeitserklärung / Declaration of Originality}

\medskip

\noindent
Hiermit bestätige ich, dass ich die vorliegende Arbeit selbständig verfasst und keine anderen als die angegebenen Hilfsmittel benutzt habe. Die Stellen der Arbeit, die dem Wortlaut oder dem Sinn nach anderen Werken (dazu zählen auch Internetquellen) entnommen sind, wurden unter Angabe der Quelle kenntlich gemacht.

\medskip

\textit{I declare that I have authored this thesis independently, that I have not used other than the declared sources / resources, and that I have explicitly marked all material which has been quoted either literally or by content from the used sources.}

\bigskip

\noindent
München, den 18.02.2024

\vspace*{2cm}

\noindent
Mehmet Aslan
}

%%% Local Variables: 
%%% mode: latex
%%% TeX-master: "d"
%%% End: 

\cleardoubleemptypage
\chapter*{Kurzfassung}
\thispagestyle{empty}

Design Patterns oder Entwurfsmuster sind in der Software-Entwicklung gängige Lösungsansätze für wiederkehrende Probleme.
Von Software-Entwicklern werden diese eingesetzt, um Probleme in der Implementierung oder in der Software-Architektur zu lösen, deren Lösungsweg bereits bekannt sind und für die jeweilige Situation abgeändert werden. Die wohl bekanntesten Design Patterns sind die 23 Entwurfsmuster nach Gamma et al.
Im weiteren Verlauf der Software-Entwicklung sind die eingesetzten Entwurfsmuster aufgrund iterativer Änderungen und steigender Komplexität im Quellcode nicht mehr einfach wiederzufinden.
Deswegen liegt der Fokus dieser Masterthesis auf der Etablierung eines Prozesses, welcher mit Machine Learning in der Lage ist, Design Pattern im Quellcode zu erkennen.

Das Ziel dieser Arbeit ist es, die Entwurfsmuster Singleton, Observer, Command und Adapter zu identifizieren. Im Kontext dieser Arbeit werden diese als Strukturen definiert, die in Summe von Rollen dargestellt werden.
Jedes dieser Rollen übernimmt im Kontext des jeweiligen Design Patterns unterschiedliche Aufgaben und Verantwortungen.
Dazu wird ein Klassifizierer trainiert, welches Rollen innerhalb dieser Entwurfsmuster durch Code-Metriken erkennt.
Anhand der Klassifikation der Rollen wird im finalen Schritt ein Übereinstimmungswert berechnet, der angibt, mit welcher Zuversicht es sich hierbei um das jeweilige Entwurfsmuster handelt.
Zum Schluss wird die Klassifikationsleistung der hier vorgestellten Methode anhand der Metriken \textit{Precision}, \textit{Recall} und \textit{f1} evaluiert. 

\bigskip

\noindent
Schlagworte: Software-Entwicklung, Machine Learning, Design Patterns 


\cleardoubleemptypage

\pagestyle{scrplain}
\pagenumbering{roman}
% ---------------------------------------------------
% D-TOC.TEX zur Verwendung mit TEXPART
% (an eigene Gegebenheiten anzupassen)
% ---------------------------------------------------
%
\tableofcontents
\clearpage
\listoffigures
\clearpage
\listoftables
\cleardoublepage


\pagestyle{scrheadings}


\addtokomafont{caption}{\small}

\mainmatter

\chapter*{Einführung}

\section*{Einführung in Design Patterns}

Bei der Entwicklung von Software-System stößt man auf Probeleme und Situationen, denen man bereits in der gleichen Form oder in einer ähnlichen Variation entgegenkommen ist.
Daher tendiert man, einen bereits bekannten Lösungsansatz hier wieder einzusetzen. Die Idee der Wiederverwendbarkeit von bekannten Lösungsanätzen für wiederaufkehrende zu lösende Probleme in der Software-Entwicklung
von Software-Systemen stellen den Kern der Entwurfsmunster oder \textit{Design Patterns} dar. Allgemeiner betrachtet bilden Entwurfsmunster eine Art Lösungsblaupause für wiederaufkehrende Probleme dar, die für die jeweilige Situation angepasst werden.
 \textit{Design Pattern - Elements of Reusable Object-Oriented Software} von Gamma et. al \cite{gamma1994design} ist wohl das bekannteste Werk, das sich mit dieser Thematik auseinander setzt.
Normalerweise werden Design Patterns verwendet, um abstrakte Probleme in der Software-Architektur von objekt-orientierten Programmiersprachen zu lösen, die sich mit der Kreation, Struktur oder Verhalten auseinandersetzen.
Obwohl Design Patterns als bekannte und erpropte Blaupause für wiederaufkehrende Probleme in Software-Systemen eingesetzt werden können, bringen Design Patterns Konsequenzen für die zu betrachtende Situation mit sich,
die in Erwägung gezogen werden sollten. Diese sind meistens erhöhte Zeit- und Speicherkomplexität als Austausch für eine bessere Flexibilität im Design des Software-Systems, so dass dessen Adaptierbarkeit für neue oder sich änderene Anforderungen sichergestellt werden können.

Die Beschreibung von Design Patterns sind detaliert und beinhalten unteranderem dessen Namen, Absicht, Motivation, Awnwendbarkeit, Struktur, Teilnehmer und deren Zusammenarbeit und andere Informationen. Die Semantik des Design Patterns inkludieren die Abischt, Motivation und Anwendbarkeit,
die angeben, was das Entwurfsmuster macht, warum dieses benötigt wird und wo es sinnvoll engegesetzt werden kann. Die Teilnehmer reflektieren deren Natur als Blaupause, da diese Rollen darstellen, die die Klassen im Kontext des Design Patterns darstellen, während die Struktur und die Zusammenarbeit zwischen den Teilnehmern die Interaktionen zwischen diesen beschreibt.
Diese Menge an Informationen ist in den Entwurfsmustern un deren Beschreibung enkodiert und wird durch Software-Entwickler in das Software-System während dessen Implementierung engeflochtet. Der Einsatz von Design Patterns ist auf Design-Entscheidung während der Entwicklung des Software-Systems zurückzuführen und des Öfteren wird diese und die jewiligen Hintergrundgedanken nicht weiter dokumentiert.
Davon ebenfalls betroffen sind Anpassungen und konkrete Implementierungsdetails des Entwurfsmusters. Schlussendlich tauchen diese Einzelheiten in Quellcode des Software-Systems ab. Das Widerfinden dieser enkodierten Information im Quellcode für die Zwecke der Weiterentwicklung und Wartung ist der Hauptmotivator für die Erkennung von Entwurfsmusterm.



\appendix
\chapter{Zusätzliche Tabellen für die Datenanalyse von P-MArt}


\begin{table}[H]
    \centering
    \begin{tabular}{|c|c|}
        \hline
        Rolle & Aufkommen in Datensatz\\
        \hline
        abstractclass & 7\\abstractfactory & 2\\abstraction & 1\\abstractproduct & 3\\adaptee & 17\\adapter & 29\\aggregate & 5\\builder & 1\\caretaker & 5\\client & 64\\colleague & 2\\command & 6\\component & 7\\composite & 11\\concreatecolleague & 4\\concreteaggregate & 6\\concretebuilder & 14\\concreteclass & 51\\concretecommand & 50\\concretecomponent & 55\\concretecreator & 13\\concretedecorator & 6\\concreteelement & 102\\concretefactory & 8\\concreteimplementor & 4\\concreteiterator & 8\\concretemediator & 2\\concreteobserver & 28\\concreteproduct & 38\\concreteprototype & 3\\concretestate & 21\\concretestrategy & 68\\concretesubject & 59\\concretevisitor & 29\\context & 24\\creator & 6\\decorator & 1\\director & 4\\element & 5\\facade & 1\\implementor & 1\\invoker & 29\\iterator & 4\\leaf & 92\\mediator & 2\\memento & 8\\nullobject & 1\\objectstructure & 4\\observer & 9\\originator & 8\\product & 46\\prototype & 1\\proxy & 5\\realsubject & 5\\receiver & 12\\refinedabstraction & 3\\singleton & 15\\state & 5\\strategy & 6\\subject & 9\\subsystemclass & 10\\target & 10\\visitor & 5\\
        \hline
    \end{tabular}
    \caption{tabellarische Darstellung der Rollenverteilung in P-MArt}
\end{table}

%%% Local Variables: 
%%% mode: latex
%%% TeX-master: "thesis.tex"
%%% End: 


\cleardoublepage

\bibliographystyle{natger}
\bibliography{thesis}

\cleardoublepage


\footnotesize
\printindex


\end{document}
